\begin{thebibliography}{Gelfand-Graev-Pyatetskii-Shapiro}
\small

\bibitem[Arthur]{Arthur} The Selberg trace formula for groups of F-rank one. \emph{Ann.\ of Math.}, to appear.

\bibitem[Arthur 2]{Arthur2} On some tempered distributions on semi-simple groups of real rank one. \emph{Ann.\ of Math.}, to appear.

\bibitem[Artin]{Artin} \"Uber eine neue Art von L-Reihen. \emph{Abh.\ Math.\ Sem.\ Univ.\ Hamburg} 3 (1923).

\bibitem[Atkin-Lehner]{AtkinLehner} Hecke operators on $\Gamma_0(m)$. \emph{Math.\ Ann.} 185 (1970), pp.\ 134-160.

\bibitem[Bargmann]{Bargmann} Irreducible representations of the Lorentz group. \emph{Ann.\ of Math.}, 48 (1947), pp.\ 568-640.

\bibitem[Borel]{Borel} Some finiteness properties of adele groups over number fields. \emph{Inst.\ Hautes \'Etudes Scient.} 16 (1963), pp.\ 5-30.

\bibitem[Cartier]{Cartier} Some numerical computations relating to automorphic functions, in \emph{Computers in Number Theory}. Academic Press, 1971.

\bibitem[Casselman]{Casselman} On some results of Atkin and Lehner. \emph{Math.\ Ann.} 201 (1973), pp.\ 301-314.

\bibitem[Cassels-Fr\"ohlichs]{CasselsFrohlichs} \emph{Algebraic Number Theory}. Thompson Book Company, Washington D.C., 1967.

\bibitem[Deligne]{Deligne} Formes modulaires et repr\'esentations l-adiques. \emph{S\'eminaire Bourbaki}, 1968/1969, n\textdegree{} 355.

\bibitem[Deligne 2]{Deligne2} La Conjecture de Weil, preprint. I.H.E.S., 1973.

\bibitem[Duflo-Labesse]{DufloLabesse} Sur la formule des traces de Selberg. \emph{Ann.\ Scient.\ \'Ecole Norm.\ Sup.}, 4$^{e}$ s\'erie, t. 4 (1971).

\bibitem[Eichler]{Eichler} Zur Zahlentheorie der Quaternionen-Algebren. \emph{Journ.\ reine angew.\ Math.} 195 (1956), pp.\ 127-151.

\bibitem[Eichler 2]{Eichler2} \emph{Lectures on Modular Correspondences}. Tata Institute, 1956.

\bibitem[Eichler 3]{Eichler3} \"Uber die Darstellbarkeit von Thetafunktionen durch Thetareihen. \emph{Journ.\ reine angew.\ Math.}, 195 (1956), pp.\ 156-171.

\bibitem[Eichler 4]{Eichler4} The basis problem for modular forms and the traces of the Hecke operators, in ``\emph{Modular Functions of One Variable I}.'' \emph{Springer Lecture Notes}, Vol. 320, 1973.

\bibitem[Gangolli]{Gangolli} Spectra of discrete subgroups, in \emph{Proceedings of Symposia in Pure Mathematics}, A.M.S. Vol. 26, pp.\ 431-436.

\bibitem[Gelbart]{Gelbart} \emph{Fourier analysis on matrix space}. \emph{Memoirs of the A.M.S.}, No. 108, 1971.

\bibitem[Gelbart 2]{Gelbart2} An example in the theory of automorphic forms, in \emph{Proceedings of Symposia in Pure Mathematics}, A.M.S., Vol. 26, pp.\ 437-439.

\bibitem[Gelbart 3]{Gelbart3} Automorphic forms on the metaplectic group, in preparation.

\bibitem[Gelbart 4]{Gelbart4} Holomorphic discrete series for the real symplectic group. \emph{Inventiones Math.} 19 (1973), pp.\ 49-58.

\bibitem[Gelfand-Fomin]{GelfandFomin} Geodesic flows on manifolds of constant negative curvature. \emph{Translations of the A.M.S.}, Series 2, Vol. I (1965), pp.\ 49-65; the Russian paper appeared in \emph{Uspekhi Mat.\ Nauk} in 1952.

\bibitem[Gelfand-Graev-Pyatetskii-Shapiro]{GelfandGraevPyatetskiiShapiro} \emph{Representation Theory and Automorphic Functions}. W. B. Saunders Co., Philadelphia, 1969.

\bibitem[Godement]{Godement} Notes on Jacquet-Langlands' theory. Institute for Advanced Study, Princeton, 1970.

\bibitem[Godement 2]{Godement2} The spectral decomposition of cusp forms, in \emph{Proceedings of Symposia in Pure Mathematics}, A.M.S., Vol. 9, pp.\ 225-234. Providence, R.I.

\bibitem[Godement 3]{Godement3} Travaux de Hecke II. \emph{S\'eminaire Bourbaki}, 1951/1952, n\textdegree{} 59.

\bibitem[Godement 4]{Godement4} Analyse spectrale des fonctions modulaires. \emph{S\'eminaire Bourbaki}, 1964/65, no. 278.

\bibitem[Godement 5]{Godement5} Les fonctions $\zeta$ des alg\`ebres simples, I, II. \emph{S\'eminaire Bourbaki}, 1958/1959.

\bibitem[Godement-Jacquet]{GodementJacquet} \emph{Zeta-functions of simple algebras}. \emph{Springer Lecture Notes}, Vol. 260, 1972.

\bibitem[Goldstein]{Goldstein} \emph{Analytic Number Theory}. Prentice-Hall, 1971.

\bibitem[Gross-Kunze]{GrossKunze} Bessel transforms and discrete series, in \emph{Springer Lecture Notes}, Vol. 276 (1972).

\bibitem[Gunning]{Gunning} \emph{Lectures on modular forms}. \emph{Ann.\ of Math.\ Studies}, Princeton Univ. Press, Princeton, 1962.

\bibitem[Harish-Chandra]{HarishChandra} Automorphic forms on a semi-simple Lie group. \emph{Proc.\ Nat.\ Acad.\ Sci., U.S.A.}, 45 (1959), 570-573.

\bibitem[Harish-Chandra 2]{HarishChandra2} Representations of semi-simple groups II. \emph{T.A.M.S.}, Vol. 76, 1954.

\bibitem[Harish-Chandra 3]{HarishChandra3} Representation theory of p-adic groups, in \emph{Proceedings of Symposia in Pure Mathematics}. Vol. 25, Providence, R.I.

\bibitem[Harish-Chandra 4]{HarishChandra4} \emph{Automorphic forms on semi-simple Lie groups}. \emph{Springer Lecture Notes}, Vol. 62, 1968.

\bibitem[Harish-Chandra 5]{HarishChandra5} Eisenstein series over finite fields, in ``\emph{Functional Analysis and related fields},'' pp.\ 76-88, Springer Verlag, 1970.

\bibitem[Hecke]{Hecke} \emph{Mathematische Werke}. G\"ottingen, 1959.

\bibitem[Howe]{Howe} Some qualitative results on the representation theory of $\GL(n)$ over a p-adic field, preprint (1972).

\bibitem[Jacquet]{Jacquet} Repr\'esentations des groupes lin\'eaires p-adiques. C.I.M.E. Summer School, Montecatini (1970).

\bibitem[Jacquet 2]{Jacquet2} \emph{Automorphic forms on GL(2), Part II}. \emph{Springer Lecture Notes}, Vol. 278, 1972.

\bibitem[Jacquet-Langlands]{JacquetLanglands} \emph{Automorphic forms on GL(2)}. \emph{Springer Lecture Notes}, Vol. 114, 1970.

\bibitem[Knapp-Stein]{KnappStein} Intertwining operators for semi-simple groups. \emph{Ann.\ of Math.} 93 (1971), pp.\ 489-578.

\bibitem[Kneser]{Kneser} Strong approximation, in \emph{Proceedings of Symposia in Pure Mathematics}, Vol. 9, pp.\ 187-196. Providence, R.I.

\bibitem[Kubota]{Kubota} \emph{Elementary Theory of Eisenstein Series}. Kodansha Ltd. and John-Wiley, 1973.

\bibitem[Kunze-Stein]{KunzeStein} Uniformly bounded representations and harmonic analysis of the $2 \times 2$ real unimodular group. \emph{Amer.\ J.\ of Math.}, Vol. 82 (1960), pp.\ 1-62.

\bibitem[Labesse]{Labesse} L-indistinguishable representations and trace formula for $\GL(2)$, preprint 1972.

\bibitem[Langlands]{Langlands} On the functional equations satisfied by Eisenstein series, mimeographed notes, Yale University.

\bibitem[Langlands 2]{Langlands2} Problems in the theory of automorphic forms, in \emph{Springer Lecture Notes}, No. 170 (1970), pp.\ 18-86.

\bibitem[Langlands 3]{Langlands3} \emph{Euler products}. Yale Univ. Press, New Haven, 1971; James Whittemore Memorial Lectures 1967.

\bibitem[Langlands 4]{Langlands4} On the functional equation of the Artin L-functions, mimeographed notes, Yale University.

\bibitem[Langlands 5]{Langlands5} Eisenstein Series, in \emph{Proceedings of Symposia in Pure Mathematics}, A.M.S., Vol. 9, pp.\ 235-257.

\bibitem[Langlands 6]{Langlands6} A little bit of number theory, notes 1973.

\bibitem[Maass]{Maass} \"Uber eine neue Art von nicht analytischen automorphen Funktionen und die Bestimmung Dirichletscher Reihen durch Funktionalgleichungen. \emph{Math.\ Ann.} 121 (1949), pp.\ 141-183.

\bibitem[Mautner]{Mautner} Spherical functions over p-adic fields. \emph{Amer.\ J.\ of Math.} Vol. 58, 1958.

\bibitem[Mautner 2]{Mautner2} Sur certaines formules de trace explicites, I-III. \emph{C.\ R.\ Acad.\ Sc.\ Paris} (1972) S\'erie A, T. 274, pp.\ 1092-1095, t. 275, pp.\ 353-356, 739-742.

\bibitem[Miyake]{Miyake} On automorphic forms on $\GL_2$ and Hecke operators. \emph{Ann.\ of Math.}, 94 (1971), pp.\ 174-189.

\bibitem[Ogg]{Ogg} \emph{Modular forms and Dirichlet series}. Benjamin, Inc., New York, 1969.

\bibitem[Robert]{Robert} Formes automorphes sur $\GL_2$ [Travaux de H. Jacquet et R. P. Langlands]. \emph{S\'eminaire Bourbaki}, 1971/1972, no. 415.

\bibitem[Roelcke]{Roelcke} \"Uber die Wellengleichungen bei Grenzkreisgruppen erster Art, \emph{Abh.\ Heidelberg Akad.\ Wiss.} 4 (1956), pp.\ 159-267.

\bibitem[Sally-Shalika]{SallyShalika} The Fourier transform on $\SL_2$ over a non-archimedean local field, preprint.

\bibitem[Sally-Warner]{SallyWarner} The Fourier transform on semi-simple Lie groups of real rank one. \emph{Acta Math.}, to appear.

\bibitem[Satake]{Satake} Spherical functions and Ramanujan conjecture, in \emph{Proceedings of Symposia in Pure Mathematics}. Vol. 9, A.M.S., Providence, R.I., pp.\ 258-264.

\bibitem[Selberg]{Selberg} Harmonic analysis and discontinuous groups in weakly symmetric Riemannian spaces with applications to Dirichlet series. \emph{J.\ Indian Math.\ Soc.}, 20 (1956), pp.\ 47-87.

\bibitem[Serre]{Serre} \emph{A Course in Arithmetic}. Springer-Verlag, 1973.

\bibitem[Shalika]{Shalika} Representations of the two by two unimodular group over local fields, thesis. Johns Hopkins University, 1966.

\bibitem[Shalika 2]{Shalika2} Some conjectures in class-field theory, in \emph{Proceedings of Symposia in Pure Mathematics}, A.M.S., Vol. 20, pp.\ 115-122, Providence, R.I.

\bibitem[Shalika-Tanaka]{ShalikaTanaka} On an explicit construction of a certain class of automorphic forms. \emph{Amer.\ J.\ of Math.}, Vol. 91 (1969), pp.\ 1049-1076.

\bibitem[Shimizu]{Shimizu} On zeta functions of quaternion algebras. \emph{Ann.\ of Math.} 81 (1965), pp.\ 166-193.

\bibitem[Shimizu 2]{Shimizu2} Theta series and automorphic forms on $\GL_2$. \emph{Jour.\ of the Math.\ Soc.\ of Jap.}, 24 (1972), pp.\ 638-683.

\bibitem[Shimura]{Shimura} \emph{Introduction to the arithmetic theory of automorphic function}. Iwanami Shoten and Princeton University Press, 1971.

\bibitem[Shimura 2]{Shimura2} On modular forms of half-integral weight. \emph{Ann.\ of Math.}, 97 (1973).

\bibitem[Shimura 3]{Shimura3} On Dirichlet series and abelian varieties attached to automorphic forms. \emph{Ann.\ of Math.}, 76 (1962).

\bibitem[Silberger]{Silberger} \emph{$\mathrm{PGL}_2$ over the p-adics}. \emph{Springer Lecture Notes}, Vol. 166, 1970.

\bibitem[Stein]{Stein} Analysis in matrix spaces and some new representations of $\SL(N,\C)$. \emph{Ann.\ of Math.}, 86 (1967).

\bibitem[Tamagawa]{Tamagawa} On $\zeta$-functions of a division algebra. \emph{Ann.\ of Math.}, 77 (1963).

\bibitem[Tanaka]{Tanaka} On irreducible unitary representations of some special linear groups of the second order, I. \emph{Osaka Journal of Math.}, Vol. 3 (1966), pp.\ 217-227.

\bibitem[Weil]{Weil} \"Uber die Bestimmung Dirichletscher Reihen durch Funktionalgleichungen. \emph{Math.\ Ann.} 168 (1967); pp.\ 149-156.

\bibitem[Weil 2]{Weil2} \emph{Dirichlet series and automorphic forms}. \emph{Springer Lecture Notes}, Vol. 189, 1971.

\bibitem[Weil 3]{Weil3} Sur certains groupes d'op\'erateurs unitaires. \emph{Acta Math.} 111 (1964), pp.\ 143-211.

\bibitem[Weil 4]{Weil4} \emph{Basic Number Theory}. Springer-Verlag, 1967.

\bibitem[Weil 5]{Weil5} Sur la th\'eorie du corps de classes. \emph{J.\ Math.\ Soc.\ Japan} 3 (1951), pp.\ 1-35.

\bibitem[Weil 6]{Weil6} Zeta functions and distributions. \emph{S\'eminaire Bourbaki}, No. 312, 1965-66.

\end{thebibliography}

\medskip
{\small \textit{Added in proof} (September 23, 1974): the reader should now also consult Volumes~II and~III of the Proceedings of the International Summer School on ``\emph{Modular Functions of One Variable},'' Springer Lecture Notes, Volumes 349 and 350, 1973; especially relevant are the papers of Casselman and Deligne in Volume~II.}
